% TODO: 英文題目と約150語の英文要旨を記入する。
\summary{Design of Dynamic Glitch Effects for Live Walpapers}{}
This study explores an interactive glitch live wallpaper for Android smartphones. Glitch expressions, such as visual noise, RGB displacement, band-like distortion, and partial image loss, were generated in real time using OpenGL ES. Instead of directly processing source images, the system uses predefined color palettes and drawing parameters to procedurally generate multiple visual layers. Device tilt is detected through onboard sensors and reflected in the wallpaper in two interaction modes: “Yura,” which moves layers at different rates to create a parallax effect, and “Zabi,” which increases noise, color displacement, flickering, and distortion in response to device movement.

A preliminary evaluation was conducted with six participants using eight wallpapers with different color, noise, parallax, and flicker settings. Each wallpaper was rated on a five-point preference scale. The results showed that preferred expressions differed among participants, suggesting that allowing users to choose among multiple glitch styles may be more appropriate than defining a single optimal expression.
